\setsrc{Źródła: [17], [14]}
\begin{frame}{Podsumowanie i zadanie praktyczne}
\begin{columns}[T,onlytextwidth]
\begin{column}{0.66\textwidth}
\codeblock{s32.py}\vspace{4pt}\consblock{s32.txt}\vspace{8pt}

{\color{plum}\bfseries Ograniczenia:}\par
\begin{itemize}\setlength\itemsep{1pt}
\item jedna osoba, 44 epoki
\item okno i pasma ustalone a priori
\end{itemize}
\end{column}
\begin{column}{0.32\textwidth}
\begin{tikzpicture}[node distance=0.25cm]
\tikzset{tk/.style={minimum width=4.4cm,align=left,text width=4.1cm}}
\node[bx,tk] (t1) {\bfseries 1.\ \texttt{subject=2}\\ {\fontsize{10}{12}\selectfont walidacja w obrębie osoby}};
\node[bx,tk,below=of t1] (t2) {\bfseries 2.\ trening S1 $\rightarrow$ test S2\\ {\fontsize{10}{12}\selectfont ocena na nowym uczestniku}};
\node[bv,tk,below=of t2] (t3) {\bfseries 3.\ porównaj wyniki\\ {\fontsize{10}{12}\selectfont parametry bez zbioru testowego}};
\end{tikzpicture}
\end{column}
\end{columns}
\end{frame}
